\section{The Religion of Europe}

\begin{quotex} I believe that the Christianization of Europe, the integration of Christianity in the European mental system was the most disastrous event in history that has happened up to this point --- a catastrophe in the true meaning of the term.
\flright{\textsc{Alain de Benoist}, \emph{The Religion of Europe}\url{https://web.archive.org/web/20120809091009/http://grece-fr.com/?p=3393}}
\end{quotex}


With that Benoist relegates 1500 years of European history as not only insignificant, but also pernicious -- and not merely pernicious, but the ``most'' pernicious event. Was it worse than the ``enlightening'' of Europe that led to the French Revolution with all its excesses and the introduction of leftism as normative? Was it worse than the Bolshevik revolution that enslaved millions in Europe for decades? Was it worse than World War II, resulting in millions of dead European men and rendering all anti-left thinking as intellectually dubious and morally reprehensible right up until our day?

If this claim is true, then the French Revolution, which was at its roots secular and anti-Christian, was a liberation of the ``European mental system'', the overturning of a catastrophe. So, we can only assume that Benoist, allegedly a ``man of the right'', would himself have ``stormed the Bastille'', shouting ``liberty, fraternity, brotherhood''.

We expect an intellectual to speak precisely, so there are no excuses to be made. It is one thing to be non-Christian for personal reasons, but quite another to make such an absurd historical claim. One suspects his intelligence is ``perverted by passion, or evil habit, or an evil disposition of nature''.

\paragraph{The Point of Support}


Benoist used to work for the \emph{Action française} as a young man and he also claims to have some 600 volumes by and about \textbf{Charles Maurras} in his personal library. So he is certainly familiar with this quote from Maurras --- also a non-Christian --- from \emph{Religious Democracy}.

\begin{quotex}
All my favorite ideas---order, tradition, discipline, hierarchy, authority, continuity, unity, work, family, corporation, decentralization, autonomy, organization of workers---had been preserved and perfected by Catholicism. As Catholicism of the Middle Ages wallowed in the philosophy of Aristotle, our social naturalism took Catholicism as one of its most solid and the dearest points of support.
\end{quotex}


Unfortunately, Benoist does not address these points, so we can only speculate. Either Benoist denies that Catholicism has historically embodied those characteristics, or else they are simply not among his favorite ideas. Specifically, if Catholic social theory was not the basis for French society, then what is? The only alternative is the secularism arising from the French Revolution, since that revolution was a pagan uprising.

So, of these qualities, which ones does Benoist reject?

\begin{itemize}


\item Order

\item Tradition

\item Discipline

\item Hierarchy

\item Authority

\item Continuity

\item Unity

\item Work

\item Family

\item Corporation

\item Subsidiarity

\item Independence

\item Organization of Workers
\end{itemize}


\paragraph{Tolerance}


Benoist seems to overlook those ideas. Like all moderns, he regards tolerance as one of the primary virtues. In his exuberance, he points to the Romans as exemplars of religious tolerance. Of course, that tolerance came at a heavy price: the ``other'' was tolerated as long as he subordinated his own gods to Rome.

\textbf{Fustel de Coulanges}, in the \textit{Ancient City}, points out that neighboring city-states were not at all tolerant of each other. Although their respective religions were superficially similar, their practice was quite different. In the event of a conflict, the victor would utterly destroy the losing city and kill all its inhabitants. So much for mutual tolerance.

Obviously, Maurras did not list tolerance as one of his favorite ideas. Nevertheless, the Catholic tradition interacted intellectually with pagans, Jews, and Muslims. They did such not out of tolerance, but rather in the recognition that they all possessed some elements of the True. That is not far from the Guenonian ideal of a fundamental unity of metaphysical principles taught by each Tradition.

\paragraph{Tellurism}


Benoist revives the old canard about being ``life affirming''. He writes:

\begin{quotex}
there is no place in paganism for a theology of \emph{exile}, based on the uprooting, the absence in the world, the absolute distance or critical negativity.
\end{quotex}


Yet \textbf{Julius Evola} describes the attitudes of men of the Solar Race:

\begin{quotex}
the type of men who naturally, through a not yet dim memory of origins and a condition of soul and body that do not cripple such a memory, feel they do not properly belong to the terrestrial race, to the extent of believing themselves men only by accident, or through ``ignorance'', or ``sleep''. The two terms vidya and avidya of the ancient Indo-Aryan teaching, that mean respectively ``knowledge'' (of the ``supreme identity'') and ``ignorance'' (that leads to the self-identification with one of the forms or modes of being of the conditioned world)
\end{quotex}


And James warns us ``to keep one's self unpolluted from this world''. (James 1:27)

Benoist doubles down on a bad bet:

\begin{quotex}
The whole Judeo-Christian theology based on the distinction of being created (the world) and uncreated being (God) ... God is realized only by and in the world.
\end{quotex}


Quite the contrary, God is realized by transcending the world. The distinction he is missing is that between ``being'' and ``becoming''. This is not just Christian, but it is likewise pagan; just read the first chapter of Evola's \textit{Revolt Against the Modern World} as described in \textit{A Priori and Ab Initio}\footnote{\url{https://gornahoor.net/?p=8408}}.

Benoist fails to grasp that in the Christian view, God is not merely transcendent, He is immanent also, since in Him ``we move and have our being''. If man is just part of the world process, as Benoist asserts, then he is unfree, the resultant of natural forces.

\paragraph{Unity}


Benoist tries to determine this about paganism:

\begin{quotex}
what \emph{inner world} it returns, what form of apprehending the world it reflects
\end{quotex}


Based on no evidence beyond his ipse dixit, Benoist claims that the ``gods'' are merely symbolic or allegorical interpretations of ``principles''. Nevertheless, they form part of the pagan mind, but they are experienced as ``external'' to him, as forces that do not arise within him but arrive, often arbitrarily, from some far off Mount Olympus. The gods may offer a boon or a bane, irrespective of his own merits. Hence, he is put in a position of trying to propitiate those gods, or forces.

On the contrary, for the Christian, man is created in the image and likeness of God. God is no arbitrary force, but the very center of his being. Moreover, God is concerned for man's well-being; there are no other gods to oppose him. Man's inner unity is not fragmented as it was for the pagan. He knows God, not as some ``wholly other'', but in knowing himself, he knows God. This self-knowledge requires the development of an inner unity, unlike the natural man, dominated by the id, who seeks out nothing beyond his bodily needs and the pleasures of the world.

\paragraph{Order}


Not surprisingly, Benoist wants to

\begin{quotex}
discard the tyranny of the \emph{Logos}, the monstrous tyranny of the Law.
\end{quotex}


The Logos, the principle of Order through which all beings have come into manifestation, is co-extensive with Being. Benoist, like all revolutionaries, finds this suffocating, a ``monstrous tyranny''. The Communist Manifesto likewise rejects the Logos, and seeks to overturn all established order.

\paragraph{Religion and Nation}


\begin{quotex}
Plato says that kinship is the community of the same domestic gods. When Demosthenes wishes to prove that two men are relatives, he shows that they practice the same religious rites, and offer the funeral repast at the same tomb. Indeed, it was the domestic religion that constituted relationship. Two men could call themselves relatives when they had the same gods, the same sacred fire, and the same funeral repast. Now, we have already observed that the right to offer sacrifices to the sacred fire was transmitted only from male to male, and that the worship of the dead was addressed to the ascendants in the male line only.
\flright{\textsc{Fustel de Coulanges}, \emph{The Ancient City}}
\end{quotex}


Without common worship, there is no kinship. Specifically, there is no Europe, and no European identity. \emph{Pace} Benoist, history bears this out. Before Charlemagne, there were only disparate and mutually incompatible tribes. Benoist's fantasy religion of Europe never existed.

\paragraph{Lost Treasures}


Benoist provides us with a litany of names that allegedly follow the ``true'' European religion. I'll mention three in particular, \textbf{Nicolas of Cusa}, \textbf{Jacob Boehme}, and \textbf{Meister Eckhart}, since they will be addressed in an upcoming post. Hence, whatever is of value in this ``religion of Europe'', sneaks into Benoist's vision because they are actually Christian value. Suddenly, in all innocence, he writes of the ``transcendental unity of the cosmos''.

He writes about Eckhart's view that God must be close to man, as if it should be a surprise to us. Benoist misrepresents Eckhart's views slightly, presumably because he does not want to have to mention Christ. That is, he claims that Eckhart asserted that ``God is born in the human soul''. Of course, what Eckhart actually asserted was that it is the Son, or Logos, that is born in the soul. That is Christian teaching, emphasized again more recently, for example, by Valentin Tomberg.

What pagan ever claimed that Zeus or Jupiter or Odin was ``born in the human soul''?

\paragraph{The Real Disaster}


\begin{quotex}
in a traditional social organization, an individual can be outside the castes in two ways, either because he is above them (ativarna) or because he is beneath them (avarna), and it is evident that these cases represent two opposite extremes. In a similar way, those among the moderns who consider themselves to be outside all religion are at the extreme opposite point from those who, having penetrated to the principial unity of all the traditions.
\flright{\textsc{Rene Guenon}}
\end{quotex}


For Tradition, the degeneration of castes is the real catastrophe. It was the overthrowing of the order of ``Throne and Altar'' that was the disaster. This means the displacement of the spiritual authority and temporal power from its rightful holders. Specifically, involves the successive overthrowing of the priestly and warrior castes in three loose stages, although they may be intermixed. The danger is that those who most vocally claim to be ``on the right'' are often part of that very degeneration

\begin{itemize}


\item \textbf{Bourgeoisie}: This caste made of farmers, merchants, and craftspeople. They can be recognized by their tendency to use the ``movement'' to sell merchandise, lectures, etc., and are interested in self-promotion.

\item \textbf{Proletariat}: When emancipated from the influence of the spiritual authority, they turn to atheism, having lost any awareness of the transcendent. Like Benoist, they become hateful of the rightful religion of Europe. Tomberg affirms that they are ``the class hostile to the hierarchical principle which is the reflection of divine order. This is why the proletariat professes atheism.'' They may go to fancy prep schools, but their mentality is prole.

\item \textbf{Outcastes}: This refers to those who are ``below the castes'', whether by birth or by lifestype. These are the opposite of Evola's well-bred man whose behavior was considered healthy and normal in a Traditional society.
\end{itemize}


The only right is the ``opposite of a revolution''. The revolution has not had the desired effect. On the contrary, it has made the proletariat even more dependent. They are kept in line with welfare programs and sedated with sex and drugs; the modern equivalent of ``bread and circuses''.

\paragraph{Conclusion}


All this goes to show that Benoist is really a Jacobin, not a man of the right. By his own admission, he does not even know what the right is, since he denies the true order of things. The ``new right'' movements always end up on the left, or at least a branch of the racist left.

Maurras returned to the faith toward the end of his life, and so, apparently, did Martin Heidegger, who is quoted extensively by Benoist. We can certainly hope for the same from Benoist and offer up a prayer.

\paragraph{Postscript on a Logical Fallacy}


It just came to our attention that at a White Privilege conference Christianity was blamed for everything bad\footnote{\url{http://dailycaller.com/2016/04/16/white-privilege-conference-almost-everything-bad-is-tied-to-christianity/}} that happened to those who are not white. It is curious that Christianity is a catastrophe for both the ``right'' and the left. Or else, it means the ``new right'' is actually the left, as we pointed out. So by that logic Christianity is a disaster for the left; the converse is also true.


\flrightit{Posted on 2016-05-02 by Cologero}

\begin{center}* * *\end{center}

\begin{footnotesize}
\begin{sffamily}
\texttt{Adam Wallace on 2016-05-02 at 05:21 said:}

Thank you for writing, Cologero; your words a light in the darkness. At some points I feel as if I'm personally being addressed based upon my own misunderstandings of now and the past, though I see that as a good thing.

\hfill

\texttt{Harun on 2016-05-02 at 08:00 said:}

The point is De Benoist is not a pagan, he's an atheist.

He may be a ``traditionalist'' but he's anti-metaphysical and thus inherently anti-Traditional (in the sense of the true Tradition).

De Benoist has only theorical knowledge of the ancient religion of Europe, he is not a practicing pagan and he even manipulate history to suit his agenda.

To him paganism is just an idea, a social construct useful to build his ideal Europe.

On these premises we understand that it has nothing to do with true Tradition as Guénon intended. This so-called paganism is not a Spiritual Tradition but an ideology for a worldy political movement.

\hfill

\texttt{Tom on 2016-05-02 at 09:47 said:}

I read Benoist's ``on being a pagan'' and was surprised that it contained nothing on that subject and was in fact a criticism of Christianity from a progressive perspective. I understand Benoist has since changed his opinion on paganism but I have not read any of his more recent works on the subject.

I agree with Harun that Benoist is not a pagan or interested in Tradition in any sense.

\hfill

\texttt{Cologero on 2016-05-02 at 13:02 said:}

Nevertheless, there is a perceived link between the so-called New Right and the Traditional authors. Benoist's flagship journal, Elements, carries articles by and about such writers. There are several web sites that try to join them, with disappointing results. Just consider some of the recent comments here.

But I hope we are not missing the larger point, since ``beating up'' on Benoist is hardly worth the trouble in itself. What I have been doing is proposing a new methodology, since the old way can hardly be understood any longer. Defenses of dogmas or moral precepts make no sense outside the proper context.

This method includes, as can be easily seen, both metaphysical and phenomenonological arguments. For that, it is necessary to move beyond merely nominal affiliations to get to the underlying worldview. Dare we call this ``radical'' traditionalism? So we point to the metaphysical elements: Being, Unity, Order or Logos, Hierarchy, and so on. So the critic of the Catholic tradition needs to identify what part of that metaphysics does he reject. As part of the phenomenonological argument, we identified the inner states of consciousness that the pagan and Catholic worldviews lead to. That is the only objective way to discuss such matters.

\hfill

\texttt{Alexander Gueric on 2016-05-03 at 14:37 said:}

The entire New Right/Alt Right scene is a joke anyways. Pseudo-intellectual cult of edgy hipster liberals pretending to be ''of the Right'' and promoting homosexuality, transexualism, atheism and larping. And as for Alain de Benoist, he is a french which is (for me) already enough to dismiss anything he says, but being a crypto Semite explains his degenerate counter-traditional agenda. Who cares what he has to say.

\hfill

\texttt{Mark Citadel on 2016-05-04 at 20:32 said:}

A very powerful argument. I have always been a little wary of the `French New Right', but this definitely puts things into perspective, and no doubt Croat Tomislav Sunic would fall under similar misconceptions as Benoist.

It is quite profound I think that Benoist would exalt `tolerance' as virtue, when we see today that the mantra of `tolerance' is the very same baseball bat swung at the head of every rightist! There are degrees of tolerance, and often it is a compartmentalized tolerance. We tolerate that which is good on measure, but the problems arise when we tolerate and worse still promote that which is bad on measure. What he maybe is exalting from Rome is not tolerance, but the Imperial Ideal, and the two ought never to be confused.

The Logos is of course a very interesting topic, especially coming as I do from a Russo mindset. In Orthodoxy, the Logos element is downplayed in favor of the unintelligible, and the opposite was the case in Roman Catholicism, yet even Roman Catholicism had Logos confined for a time to its proper station. There is, in such discussions, too much absolutism and not enough particularism. Destroy Logos? Absurd. Even those theories which propose chaos as a weapon and a rhetorical tool are clear to say that the objective is to save Logos from itself, to save the logical and analytical view of an intelligible world from the crushing vice it has designed for itself in the Age of Reason!

Pagan civilizations of all stripes had some degree of the spirit of Logos, for had they not, they never would have advanced beyond subsistence level. The great Pagan empires would never have existed without that spirit. Man cannot live on mysticism alone, and even I, a lover of all things mystical, can admit that. We need a reconciliation of the intelligible with the unintelligible, that which can be studied and that which must be shrouded in secrecy. ``Everything in its proper place'' is the clarion call of the Reactionary.

On a side note, concerning the nature of Europe, I was quite interested to learn about the origins of the term. The goddess Europa was Phoenician, and originally European was to denote all those surrounding the Mediterranean Sea. It would not have included Slavs or Nordics etc. and probably not even Benoist himself. It seems religion and geopolitics have shaped `Europe' more than race has. The silly idealizations of Europeans as a bloc have to be abandoned. If we speak of an Occidental, or even a Eurasian civilization, then we are speaking of empire and its attendant theopolitical factors. The Kingdom of God itself is not built without Logos. Why do we think we could fashion earthly kingdoms without it? As above, so below.

\hfill

\texttt{nickbsteves on 2016-05-05 at 12:18 said:}

\begin{quote}\footnotesize

If this claim is true, then the French Revolution, which was at its roots secular and anti-Christian, was a liberation of the ``European mental system'', the overturning of a catastrophe. So, we can only assume that Benoist, allegedly a ``man of the right'', would himself have ``stormed the Bastille'', shouting ``liberty, fraternity, brotherhood''.
\end{quote}


Not necessarily. One can oppose X \emph{and} Y, even if they are hostile to each other, and most especially if one thinks that X was somehow responsible for Y. I don't know how much blame Benoist heaps upon Christianity for the European Revolutions. But if much, he wouldn't be the first. I would attach much blame for them upon the Protestant Reformation, which is to say a heretical version of Christianity. How much blame does Christianity deserve for its particular heresies? I don't actually know. But I'm not convinced the question really matters any more than assigning blame for the lethality of fire ants to natural selection. What's done is done. Europe is Christian, for better or worse. And there's no putting the genies back in their respective bottles.

\hfill

\texttt{Cologero on 2016-05-05 at 12:56 said:}

The point, nickbsteves, is that if X is the worst evil, then, even if Y is also evil, Y would be preferable to X on the principle of ``the lesser of two evils''. That is necessarily so. Even if you argue that ``X was somehow responsible for Y'' (its polar opposite, by the way), Y is still an improvement. The issue has nothing to do with ``assigning blame'' but rather in the ordering of disasters as in \textit{Lattice Theory}\footnote{\url{https://en.wikipedia.org/wiki/Lattice_(order)}}.


The issue is important for the reasons I pointed out (which are far from what apparently concerns you). It gets to the root of what the ``right'' is. If it is fundamentally about a Traditional social order, as I asserted, as well as a certain metaphysical viewpoint, these are then part of the worst disaster in European history as B claims. So I'll ask you the same question I posed to Mr. B: which do you reject? I don't think what I wrote was all that difficult to understand.

\hfill

\texttt{Nilakantha on 2016-05-05 at 14:33 said:}

I generally use the categories of philosophy, religion, art and science to judge the degree to which any society is civilized. And to show my cards, I equate orthodox philosophy with the developed Platonism that begins either with Plotinus or his teacher Ammonius Saccas. The fall into barbarism is a process by which a social order falls into the unintelligibility of the aforementioned categories. In Europe I see this fall beginning with the move from Augustinian to Thomistic theology; it was accelerated by the Renaissance and Reformation, and became unstoppable in the Enlightenment.

A return to civilization in Europe is only possible with a radical turning to orthodox philosophy and orthodox religion. In the case of Europe this means a return to Catholic Christianity (Catholic here means Roman Catholic, as well as Eastern and Oriental Orthodoxy). Since European paganisms long ago lost any effective connection to metaphysical reality, any attempt to base a renewed Europe on pagan practice isn't just wrong; it's wrongheaded. Since to be civilized means to order a culture according to transcendent principles, an orthodox religion must be in place so that transcendent principles may be effectively transmitted to the psychic and material realms. Only Catholic Christianity can play this role for Europe.

\hfill

\texttt{traxler on 2016-05-06 at 07:58 said:}

Europe's future religion will need to be some form of Vedicism. Any religion which denies, or ignores karma, dharma, and Reincarnation is hollow and incomplete.

The Vedic teachings corroborate the remnants of ancient European Paganism still left with us today. Both are native to the Aryan spirit and in harmony with Nature.

Judaic religions, such as Christianity, are alien to the Aryan spirit. These false religions with their fallacious univeralism lead to rebellion against nature. It starts with meddeling in Africa (2 billion negroid pests thanks to Catholicism), and it ends with our annihilation, as you see today in Europe. Liberalism is not possible without the Christian axiology upon which it is based.

There was nothing traditional or harmonious about Europe's bloody transition into Christianity. Christianity, like Bolshevism, was violently revolutionary and apocalyptic. It conquered by the sword.

It cut down Europe's living symbol of the tree (Yggdrasil) and replaced it with the cross--a symbol of torture, horror, and death. As Hitler stated in Table Talks, and George Orwell parallels in 1984: Christianity was the first to exterminate its enemies in the name of love.

The Bible was mostly written by Jews, why anyone would ever make the argument that this is a suitable theology to hedge Europe's survival on is beyond me. Reason doesn't figure into it as plenty of otherwise intelligent Christians I've known abandon all reason when confronted by the above.

You will never win against the JJew until you remove him from your soul first. Christianity is spiritual cuckoldery.

\hfill

\texttt{Cologero on 2016-05-06 at 17:44 said:}

Traxler:

We have written about karma and dharma in \textit{Predestination and Predilection}\footnote{\url{http://www.gornahoor.net/?p=3998}}, so not a problem.

I think I'm beginning to understand you. You want to worship monkeys and rats, put a statue of Ganesh on the dashboard of your car, and be vigilant to avoid letting a Jewish thought insinuate itself into your consciousness. Not just that ... then you want to persuade your friends, family, and eventually all the peoples of Europe to adopt your viewpoint. Keep us posted every 6 months or so about your progress.

Your moralizing tone about past events is very sweet, and would endear you into the liberal camp. The Greeks conquered by force, as did the Romans, Vikings, and Germanics. So, too, did Constantine, Charlemagne and Clovis. Read some Machiavelli and learn how men of power behave. The enemy either capitulates or is destroyed. Tolerance is a modern invention. I suspect you are more influenced by Marcuse than you realize.

\hfill

\texttt{Mark Citadel on 2016-05-06 at 17:46 said:}

@traxler -- Good to see Uncle Adolf is still table-talking. Will you people ever learn? I'm sure your \#1 article this year was

`Charlemagne: Defender of Europe, or Servant of the JOOOS!'

\hfill

\texttt{traxler on 2016-05-06 at 20:41 said:}

Even worshipping monkeys and rats is superior to worshipping a jew.

No I have no problem with force, why else would I quote Hitler if I were a pacifist.

What I do have a problem with is Jewish thought penetrating the Aryan conscience through a Jewish religion.

Why did you resort to ridicule instead of addressing this main point?

\hfill

\texttt{Me on 2016-05-06 at 20:58 said:}

I have read a lot of Jewish thoughts, but rarely are they morally universalist

\hfill

\texttt{traxler on 2016-05-07 at 11:04 said:}

You mustn't be paying much attention then. In case you didn't realize it the moral universalism of the New Testament, Marxism, Liberalism, Feminism, and anti-racism was all authored by Jewish hands.

Though I understand the above explanation is lost on someone like you with such poor reading comprehension. But the above still stands for the benefit of others, not you.

\hfill

\texttt{Me on 2016-05-07 at 12:21 said:}

My response embarrasses me, since it is ``low hanging fruit'' and the question of Christianity's ``Jewishness'' is the most vague, irrelevant and uninteresting point in your comments. Yet you called it your main point, so I probably shouldn't be too hard on myself. I read your real points as:

1. Christianity is a universalist religion which leads to liberalism.

2. Christianity is not ``native to the Aryan spirit and in harmony with Nature''.

The second objection is too vague for me, though you didn't ask me personally to address it. As for the first point, since I am just a nobody memer, you'd probably be more interested in reading Cologero's thoughts: \url{http://www.gornahoor.net/?p=265}

\hfill

\texttt{Cologero on 2016-05-14 at 22:54 said:}

Unfortunately, traxler, reading comprehension is not of much use unless one actually reads something. So this is for the benefit of others perhaps like you --- and there are many --- who may read this and become ashamed at their lack of knowledge. First of all, the historian Ernest Renan, whose intelligence and erudition are not in doubt, tells us a quite different story.

\begin{quote}\footnotesize

It was then that the religious ideas of the races grouped around the Mediterranean became profoundly modified; that the Eastern religions everywhere took precedence; that the Christian Church, having become very numerous, totally forgot its dreams of a millennium, broke its last ties with Judaism, and entered completely into the Greek and Roman world.
\end{quote}


To lump liberalism with the Tradition is quite ignorant since the roots of their worldviews are incompatible. The difference, for example, between a St Augustine and a Rousseau, is unbridgeable. To hear from Edmund Burke, who was a witness to the French Revolution, understood its foundations. He wrote:

\begin{quote}\footnotesize

Instead of the religion and the law by which they were in a great politick communion with the Christian world, they have constructed their Republick on three bases, all fundamentally opposite to those on which the communities of Europe are built. Its foundation is laid in Regicide; in Jacobinism i.e., liberty, equality, fraternity ; and in Atheism; and it has joined to those principles, a body of systematick manners which secures their operation.
\end{quote}


Ananda Coomaraswamy explains that it is hardly likely that a European who does not even understand his own tradition would be able to understand the Vedanta:

\begin{quote}\footnotesize

The educated man of today is completely out of touch with those European modes of thought and those intellectual aspects of the Christian doctrine which are nearest those of the Vedic traditions. A knowledge of modern Christianity will be of little use because the fundamental sentimentality of our times has diminished what was once an intellectual doctrine to a mere morality that can hardly be distinguished from a pragmatic humanism. A European can hardly be said to be adequately prepared for the study of the Vedanta unless he has acquired some knowledge and understanding of at least \emph{Plato, Philo, Hermes Trismegistus, Plotinus, Gospel of John, Dionysius the Areopagiite, Meister Eckhart, Dante}. Eckhart, with the possible exception of Dante, can be regarded from an Indian point of view as the greatest of all Europeans.
\end{quote}


So how many of those works are you familiar with? Has not Gornahoor dealt with all of them, as well as with the Vedanta. That is a conversation you could join if you had the requisite background.

In another case, the French philosopher Jean Borella, began the study of Rene Guenon as a young man, which led to his study of the Vedanta. However, this simply led him to appreciate his own tradition all the more. He wrote about his experience:

\begin{quote}\footnotesize

I went back to ancient doctrines like a delighted child going from discovery to discovery, from treasure to treasure, from marvel to marvel. I recognized and loved this Christian past, its beauty no unworthy of the God whom it had honored with its liturgy, cathedrals and theologies. It was in me as flesh of my flesh, soul of my soul, hear of my heart, and I did not know it. Once discovered, fixing the gaze of my spirit upon the holy Fathers and Doctors, upon the Clements, the Dionysii, the Gregories, the Augustines, and the Thomases, I said: \textbf{I too am of their race. Surely not by sanctity or genius, but by blood.} Drinking in the freshness of the ages, I felt my Christian soul revive.
\end{quote}

\hfill

\end{sffamily}
\end{footnotesize}

